% ECONOMICS_PROBLEM: {"id": "D83-571", "title": "Human-aware model training", "jel_code": "D83", "source_id": "OP-0418"}
% !TeX program = xelatex
\documentclass[11pt,a4paper]{article}
\usepackage[margin=23mm]{geometry}
\usepackage{fontspec}
\setmainfont{Latin Modern Roman}
\usepackage{amsmath,amssymb,mathtools}
\usepackage{xcolor,enumitem,booktabs,longtable,array,xurl,hyperref}
\definecolor{muted}{HTML}{53636C}
\hypersetup{colorlinks=true,urlcolor=blue,linkcolor=blue}
\setlength{\parindent}{0pt}
\setlength{\parskip}{5pt}
\newcommand{\R}{\mathbb R}
\newcommand{\E}{\mathbb E}
\newcommand{\N}{\mathbb N}
\newcommand{\ind}{\mathbf 1}
\newcommand{\Var}{\operatorname{Var}}
\newcommand{\Cov}{\operatorname{Cov}}
\newcommand{\supp}{\operatorname{supp}}
\DeclareMathOperator*{\argmin}{arg\,min}
\DeclareMathOperator*{\argmax}{arg\,max}
\newcommand{\entrymeta}[3]{{\sffamily\small\color{muted}\textbf{\texttt{#1}}\quad|\quad #2\par #3\par}}
\newcommand{\Problemref}[1]{\texttt{#1}}
\newcommand{\JELref}[2]{\texttt{#2} (JEL \texttt{#1})}
\begin{document}
\section*{Human-aware model training}
\textbf{Problem ID:} \texttt{D83-571}\quad \textbf{JEL:} \texttt{D83}\par
% BEGIN ECONOMICS STATEMENT
\label{problem:OP-0418}
{\sffamily\small\color{muted}Original subject group: Innovation and AI economics\par}
\label{definitions:problem:30007700}
\textbf{Audit classification:} Resolved or misstated baseline. The paper exists, with Kevin A. Bryan as coauthor. The original proposed heterogeneity extension overlaps analyses already in Sections 4 and 5; only the additional robust/identification specification is new and its open status is unverified.
\subsection*{Definitions and precise research target}
Let economic states \(\omega\) follow a specified distribution \(P\), and let a model trained with parameter \(\theta\in\Theta\) generate signal \(S_\theta\) according to a known signal experiment. Users have types \(v\) drawn from distribution \(F\), where type determines verification cost \(k_v\), available corrective information and action set \(\mathcal A_v\). Let \(u_v(a,\omega)\) be a specified bounded economic payoff. After observing the signal, a user may verify, override or accept the suggested action. Let \(a_v^\ast\) and \(e_v^\ast\in\{0,1\}\) be the realized action and verification indicator under an optimal contingent response policy after observing \(S_\theta\), with actions allowed to depend on verification results and ties resolved by a stated rule. Training incurs real cost \(C(\theta)\). The proposed training objective is
\[
\max_{\theta\in\Theta}E_Q[u_v(a_v^\ast,\omega)-k_v e_v^\ast]-C(\theta).
\]
The joint law \(Q\) has the specified state and user-type marginals; all expectations integrate states, model randomness and user types. Extend this objective to a declared class \(\mathcal Q\) of plausible state and user distributions where \(\mathcal Q\) is a given finite set of candidate distributions, by maximizing its minimum over \(\mathcal Q\), subject to a stated harm constraint. Determine when economically optimal training differs from minimizing average prediction error, and which user-response data identify the relevant payoff and verification functions. Existing work already solves stylized human-aware training problems; heterogeneous users are also analyzed in that work; the additional distribution-robust identification specification here is an editorial extension whose unresolved status is not independently verified.

\textbf{Definitions, assumptions and mathematical qualifications.} Let finite state, user-type and action sets have joint law \(Q\); independence of state and type is imposed only when desired. A training choice \(\theta\) specifies signal kernel \(K_\theta(s\mid\omega,v)\) and verification kernel \(K_v^{\rm ver}(h\mid\omega,s)\). A user policy maps its observed type and signal to verification and then maps the available information to an action. Define \(V_Q(\theta)\) as maximum expected bounded payoff less verification cost over these policies. For a nonempty finite uncertainty set \(\mathcal Q\), robust training is \(\sup_{\theta\in\Theta}\min_{Q\in\mathcal Q}\{V_Q(\theta)-C(\theta)\}\), with \(E_Q[h(a,\omega)]\leq\bar h\) for all \(Q\). Finite \(\Theta\), or compact \(\Theta\) with suitable continuity, is required for an attained maximum. Specify whether users know \(Q\), learn it, or optimize their own worst-case criterion; those produce different \(V_Q\).
\subsection*{Variants and assumption changes}
\begin{enumerate}
\item \textbf{Sourced solved baseline.} Bryan and Gans analyze a coverage--accuracy frontier and human continuation decisions, including heterogeneous agents and teams. These ingredients alone must not be relabelled unresolved.
\item \textbf{Editorial extension.} Unknown response primitives are estimated from randomized signal/verification data; optimize distribution-robust training with a uniform harm constraint and quantify policy regret over compatible user models.
\end{enumerate}
\subsection*{References and source audit}
\textbf{Checked source.} arXiv:2608.12538v1, submitted 12 August 2026, Sections 2--5 and 6. The paper analyzes optimal training, heterogeneous users, verification, monopoly and teams. Full primary HTML retrieved. \url{https://arxiv.org/html/2608.12538v1}.
\begin{enumerate}\raggedright
\item\hypertarget{definitions:source:30007700:1}{} Kevin A. Bryan; Joshua S. Gans. 2026. \emph{Training AI for When Humans Will Use It}. arXiv:2608.12538v1, submitted 12 August 2026, Sections 2--5 and 6.
\url{https://arxiv.org/html/2608.12538v1}.
DOI: \url{https://doi.org/10.48550/arXiv.2608.12538}.
\end{enumerate}

\subsection*{Common conventions: Source mathematical conventions}

These conventions apply to each entry unless explicitly replaced.

\textbf{Models and identification.} A primitive model \(m\) specifies all probability laws, feasible choices, information, equilibrium and measurement rules used by its target. The admissible class \(\mathcal M\) is fixed independently of outcomes. If \(P_O(m)\) is its observable probability law and \(\tau(m)\) the target, its identified set at observable law \(P\) is
\[
 \mathcal I_\tau(P)=\{\tau(m):m\in\mathcal M,\ P_O(m)=P\}.
\]
Point identification means that this set is a singleton. An empty set signals incompatibility of data and assumptions. Coordinate bounds need not describe the joint set. Bounds are sharp only if every retained target is attainable or arbitrarily approachable by compatible models. Finite bounds require explicit restrictions on unobserved quantities; unrestricted confounding, hidden wealth, extrapolation or arbitrary equilibrium selection can yield uninformative sets.

\textbf{Causal interventions.} A potential outcome \(Y_i(p)\) is the outcome under a complete policy regime \(p\), including timing, financing, information and market spillovers. The target population and units are fixed before treatment. With interference, use the complete assignment vector or a justified exposure mapping; individual no-interference notation alone cannot identify an economy-wide intervention. A difference of mean potential outcomes does not require the cross-world joint distribution. Conditioning on post-treatment migration, survival, enrollment or employment changes the estimand and needs separate justification.

\textbf{Statistical statements.} An estimator, confidence set, forecast or test specifies a sampling law, model class and loss/coverage criterion. Uniform \((1-\alpha)\)-coverage means \(\inf_{m\in\mathcal M}\Pr_m\{\tau(m)\in C_n\}\geq1-\alpha\), or an explicitly asymptotic version. The whole parameter space trivially has coverage, so informativeness requires a stated width, risk or power objective. A forecasting improvement is relative to a named benchmark, information set and evaluation population.

\textbf{Equilibrium and optimization.} An equilibrium comprises feasible plans, optimality, beliefs where needed, budget constraints and all market/resource clearing. The primitives or a stated benchmark define each correspondence \(\mathcal E(p;m)\). Nonemptiness is assumed or proved, never inferred just from compact policy sets. A selected-equilibrium target uses a declared selection rule; a robust target quantifies over all permitted equilibria. Use suprema or infima unless attainment is established. An \(\varepsilon\)-optimal policy is feasible and within \(\varepsilon\) of the finite optimum value. Report an empty feasible set separately.

\textbf{Welfare and accounting.} Normative utility, interpersonal normalization, social weights, numeraire, discounting and the use of public receipts are primitives. Behavioral utility need not coincide with the welfare criterion. Household welfare includes ownership income and rents once; adding profits again to compensating variation generally double-counts them. A monetary equivalent is defined by an explicit transfer and price convention, with monotonicity and range conditions. Average effects, mechanism contributions, transfers and welfare losses are different targets. Infinite sums require convergence and dynamic feasible plans require terminal or no-Ponzi conditions.

\textbf{Source versions.} A working-paper date, revision date and journal publication year are distinguished. A DOI for a journal version is not automatically the DOI of an earlier manuscript. Locators refer to the stated version and printed pages where specified. A metadata-only check does not verify a claim about a full-text conclusion. Every entry's source subsection narrows the provenance claim accordingly.
% END ECONOMICS STATEMENT
\end{document}
